\noindent\textbf{(a)} Debe existir una función continua $h(v)$ tal que $P_y-Q_x=h(x^3+y^3)(3x^2Q-3y^2P)$. Entonces $\mu=A\exp\!\int h(v)\,dv$, $A\ne0$. Donde el denominador no se anula, equivale a que $(P_y-Q_x)/(3x^2Q-3y^2P)$ sea una función de $v$.

\noindent\textbf{(b)} $\mu_1=1/(x^3+y^3)$ en $x+y\ne0$. Una integral primera es $F=(x^2-xy+y^2)/(x+y)=C$. Se añade la solución $y=-x$, excluida por el factor.

\noindent\textbf{(c)} El criterio es $P_y-Q_x=h(x+y)(Q-P)$. Para la aplicación $h(v)=-2/v$ y $\mu_2=1/(x+y)^2$.

\noindent\textbf{(d)} $\mu_2/\mu_1=(x^2-xy+y^2)/(x+y)=C$, junto con $y=-x$ y las prolongaciones diferenciables verificadas en la forma original.
